\section{Background}

\subsection{Sequential Bayesian Continual Learning}

Consider a sequence of $T$ non-overlapping tasks with i.i.d datasets $\mathcal{D}_{1:T}=(\mathcal{D}_1,\ldots,\mathcal{D}_T)$ and a model with parameters $\theta\in\mathbb{R}^P$. We assume that the same parameter space is used throughout the task sequence and after learning task $t$, its data are not available when training future tasks. Let $q_0(\theta)=p(\theta)$ be the initial prior and let $q_t(\theta) = p(\theta\mid\mathcal{D}_{1:t})$ denote the posterior after task $t$. If the task datasets are conditionally independent given $\theta$, Bayesian updating gives
\begin{equation}
    q_{t+1}(\theta)
    = \frac{p(\mathcal{D}_{t+1}\mid\theta)\,q_t(\theta)}
           {\int p(\mathcal{D}_{t+1}\mid\theta')\,q_t(\theta')\,d\theta'}.
    \label{eq:sequential_bayes}
\end{equation}
Thus, under the assumed model, $q_t$ contains the information from earlier tasks needed for the next Bayesian update.

For most applications, storing past tasks' samples is infeasible due to privacy or storage concerns and evaluating $q_t$ exactly is intractable. A replay-free learner instead retains a tractable approximation $\hat q_t$ after discarding $\mathcal{D}_t$. The new task then defines a target proportional to $p(\mathcal{D}_{t+1}\mid\theta)\hat q_t(\theta)$, but this expression does not directly provide samples from the target distribution. Moreover, the
approximation becomes the prior for the following task, so errors in $\hat q_t$ can propagate through subsequent updates~\cite{kessler2023sequential}. Probabilistic CL must therefore both represent the accumulated posterior and update that representation as new tasks arrive.

\subsection{Schr\"odinger Bridges}

Let $s\in[0,1]$ and $t\in\mathbb{Z}^+$ denote the time index of a stochastic process on $\mathbb{R}^P$ and the task index, respectively. Given two marginal distributions
$\pi_0$ and $\pi_1$ and a reference process $\mathbb{Q}$, the SB is the path measure closest to $\mathbb{Q}$ that connects the two
marginals~\cite{leonard2014survey}:
\begin{equation}
    \mathbb{P}^\star = \arg\min_{\mathbb{P}}\ \mathrm{KL}(\mathbb{P}\,\|\,\mathbb{Q})
    \quad \text{s.t.} \quad \mathbb{P}_0=\pi_0,\ \ \mathbb{P}_1=\pi_1 .
    \label{eq:sb}
\end{equation}
We consider references given by a diffusion
$d\theta_s = b(\theta_s,s)\,ds + \sigma\,dW_s$, where $b$ is the reference drift
and $W_s$ a standard Wiener process. If $\mathbb{P}$ has drift $b+u$ with the
same $\sigma$, Girsanov's theorem gives
\begin{equation}
    \mathrm{KL}(\mathbb{P}\,\|\,\mathbb{Q})
    = \mathrm{KL}(\mathbb{P}_0\,\|\,\mathbb{Q}_0)
    + \frac{1}{2\sigma^2}\,\mathbb{E}_{\mathbb{P}}\!\int_0^1 \|u(\theta_s,s)\|^2\,ds .
    \label{eq:girsanov}
\end{equation}
Since the first term is fixed by $\pi_0$, the SB is the process that reaches
$\pi_1$ with the least control effort added to the reference dynamics. Its drift
takes the form $b+\sigma^2\nabla_\theta\log\varphi_s$, where the potential
$\varphi_s$ is determined by the two marginals.

\section{Posterior Schr\"odinger Bridges}
\label{sec:method}

PSB represents the posterior after each task by a Gaussian mixture centered on
an ensemble of networks. When task $t+1$ arrives, annealed Langevin dynamics
moves this population toward the updated target, guided by the bridge learned at
the previous boundary. PSB then fits a new bridge between $\hat q_t$ and the
updated population under a Fisher-metric Brownian reference, and the updated
population defines $\hat q_{t+1}$.

\subsection{Posterior representation}
\label{sec:representation}

We approximate the posterior after task $t$ by a mixture of $K$ Gaussian
components centered on an ensemble of networks, which we call particles,
\begin{equation}
    \hat q_t(\theta)=\sum_{k=1}^{K}w_t^{(k)}\,
    \mathcal{N}\!\left(\theta;\,\theta_t^{(k)},\,A_t^{-1}\right).
    \label{eq:mixture}
\end{equation}
The components allow the approximation to cover distinct parameter solutions,
and the shared diagonal precision $A_t$ constrains changes to weights that are
important for earlier tasks. For the first task, the particles are initialized
independently from a He-uniform distribution~\cite{he2015delving}, trained on
$\mathcal{D}_1$ with their own sequences of minibatches, and treated as equally
weighted samples of $q_1$, following the view of deep ensembles as approximate
Bayesian model
averaging~\cite{lakshminarayanan2017ensembles,wilson2020bayesian}. After every
task, the precision accumulates the Fisher information of that task,
\begin{equation}
    A_t=\alpha I+\sum_{j=1}^{t}\gamma^{\,t-j}N_jF_j,
    \label{eq:precision}
\end{equation}
where $N_j=|\mathcal{D}_j|$ is the number of examples in task $j$ and
$F_j=\mathbb{E}_{(x,y)\in\mathcal{D}_j}\,\mathbb{E}_{k}\big[\big(\nabla_\theta\log
p(y\mid x,\theta_j^{(k)})\big)^{\odot2}\big]$ is its diagonal empirical Fisher
averaged over particles, so that $N_jF_j$ is the Laplace precision of the
likelihood of task $j$. Here $\gamma\in(0,1]$ discounts older tasks, with
$\gamma=1$ giving an undiscounted Laplace propagation, and $\alpha>0$ is the
precision of a Gaussian prior that keeps $A_t$ invertible.

\subsection{Target posterior and its score}
\label{sec:target}

The target posterior is the Bayesian update of the carried posterior,
\begin{equation}
    \tilde q_{t+1}(\theta)\propto p(\mathcal{D}_{t+1}\mid\theta)\,\hat q_t(\theta).
    \label{eq:target}
\end{equation}
Applying the full likelihood of the new task in a single step
would concentrate the particle weights on a few networks and require the
particles to travel the entire distance to the target at once. We therefore
introduce the likelihood gradually through the tempered posteriors
\begin{equation}
    \pi_{t,\beta}(\theta)\propto p(\mathcal{D}_{t+1}\mid\theta)^{\beta}\,
    \hat q_t(\theta), \qquad \beta\in[0,1],
    \label{eq:tempered}
\end{equation}
which begin at $\pi_{t,0}=\hat q_t$ and end at $\pi_{t,1}=\tilde q_{t+1}$, so that
$\beta=1$ recovers the Bayesian update exactly and each small increase of $\beta$
changes the distribution only slightly. The score of the tempered posteriors
separates into the contributions of the new task and of the carried posterior,
\begin{equation}
\begin{split}
    \nabla_\theta\log\pi_{t,\beta}(\theta)
    &= \beta\,\nabla_\theta\log p(\mathcal{D}_{t+1}\mid\theta)\\
    &\quad -A_t\Big(\theta-\textstyle\sum_{k=1}^{K}r_k(\theta)\,\theta_t^{(k)}\Big),
\end{split}
\label{eq:pathscore}
\end{equation}
where $r_k(\theta)\propto w_t^{(k)}\mathcal{N}(\theta;\theta_t^{(k)},A_t^{-1})$
is the probability that $\theta$ belongs to component $k$. The first term is the
likelihood score, which drives the particles toward the incoming task. The
second is the prior score, which pulls each particle toward the carried
solutions it most likely belongs to, with a strength set by the accumulated
Fisher information.

For a single component, the log-target is
$\log p(\mathcal{D}_{t+1}\mid\theta)-\tfrac12\|\theta-\theta_t\|_{A_t}^2$ up to a
constant, with $\|u\|_M^2=u^{\top}Mu$. Substituting (\ref{eq:precision}),
negating and dividing by $N_{t+1}$, maximizing the log-target is equivalent to
minimizing
\begin{equation}
    \bar\ell_{t+1}(\theta)
    +\tfrac{1}{2}\textstyle\sum_{j=1}^{t}\gamma^{\,t-j}\tfrac{N_j}{N_{t+1}}
    \|\theta-\theta_t\|_{F_j}^{2},
    \label{eq:ewc}
\end{equation}
where $\bar\ell_{t+1}$ is the mean negative log-likelihood of the new task and
the small isotropic term $\tfrac{\alpha}{2N_{t+1}}\|\theta-\theta_t\|^2$ is
omitted. This is the online EWC
objective~\cite{kirkpatrick2017ewc,schwarz2018progress} with penalty strength
fixed at $N_j/N_{t+1}$, so the EWC solution is the most probable parameter under
the target of PSB with a single particle.

\subsection{Annealed Langevin transport}
\label{sec:langevin}

As $\beta$ increases from $\beta_{m-1}$ to $\beta_m$, each particle weight is
multiplied by the change in the tempered likelihood,
$p(\mathcal{D}_{t+1}\mid\theta^{(k)})^{\beta_m-\beta_{m-1}}$, with each increment
as large as possible while the effective number of particles
$1/\sum_k(\bar w^{(k)})^2$ of the normalized weights stays above a fixed
fraction of $K$. When the weights become uneven, low-weight particles are
replaced by copies of high-weight ones and the weights are reset. The particles
then take preconditioned Langevin steps toward $\pi_{t,\beta_m}$,
\begin{equation}
    \theta'=\theta+\eta\,A_t^{-1}\nabla_\theta\log\pi_{t,\beta_m}(\theta)
    +\sqrt{2\eta}\,A_t^{-1/2}\xi,
    \label{eq:mala}
\end{equation}
with $\xi\sim\mathcal{N}(0,I)$ and an adapted step size $\eta$, each accepted
with a probability computed from $\pi_{t,\beta_m}$ on the full task data. The
drift of these dynamics is the score (\ref{eq:pathscore}), and the
preconditioner makes the steps cautious in directions constrained by earlier
tasks. At $\beta=1$, the particles and their normalized weights form the
terminal population $\{\Theta_1^{(k)}\}$.

\subsection{Fisher-metric Schr\"odinger bridge}
\label{sec:bridge}

At each boundary, PSB solves (\ref{eq:sb}) with $\pi_0=\hat q_t$,
$\pi_1=\tilde q_{t+1}$ and the Fisher-metric Brownian reference
\begin{equation}
    \mathbb{Q}_t:\quad d\Theta_s=\sigma A_t^{-1/2}\,dW_s,\qquad s\in[0,1],
    \label{eq:fisherreference}
\end{equation}
where $s$ is the bridge time, running from $\hat q_t$ at $s=0$ to
$\tilde q_{t+1}$ at $s=1$, and $\sigma>0$ sets the noise level. We denote its
solution by
$\mathbb{P}^\star_t$. With this noise, the control cost in (\ref{eq:girsanov})
becomes $(2\sigma^2)^{-1}\mathbb{E}\int_0^1\|u_s\|_{A_t}^2\,ds$. The reference
therefore assigns a higher cost to movement in directions that earlier tasks
constrain strongly, so that the stability--plasticity trade-off enters as a
transport cost. Solving this problem at every task boundary forms the chain of
SBs.

We solve each bridge by iterative Markovian fitting
(IMF)~\cite{shi2023dsbm,peluchetti2023idbm}. Given a coupling $\Pi$ of samples
from $\hat q_t$ and the terminal population, the reciprocal step draws
intermediate parameters from the conditional Brownian bridge of $\mathbb{Q}_t$,
\begin{equation}
    \Theta_s=(1-s)\,\Theta_0+s\,\Theta_1+\sigma\sqrt{s(1-s)}\,A_t^{-1/2}\xi,
    \label{eq:conditionalbridge}
\end{equation}
with $\xi\sim\mathcal{N}(0,I)$, and the Markovian step fits the forward drift by
\begin{equation}
    \min_\phi\ \mathbb{E}\left\|v^{\mathrm{f}}_\phi(\Theta_s,s)
    -\frac{\Theta_1-\Theta_s}{1-s}\right\|_{A_t}^{2},
    \label{eq:lossf}
\end{equation}
where $(\Theta_0,\Theta_1)\sim\Pi$ and $s$ is drawn uniformly from $[0,1]$. The
backward drift $v^{\mathrm{b}}_\psi$ is fitted in the same way with target
$(\Theta_0-\Theta_s)/s$, and the $A_t$-norm matches the control cost of the
reference. A new coupling is obtained by simulating the fitted forward process
from samples of $\hat q_t$, or the fitted backward process from the terminal
population, alternating directions between rounds. The coupling is initialized
with independently drawn endpoint samples. Exact IMF projections preserve the
marginals and converge to $\mathbb{P}^\star_t$, whereas learned drifts
approximate them. We therefore measure the energy distance between the
simulated endpoint and the terminal population.

A drift over all weights at once would be too large to fit from $K$ particles,
so the drifts act on neuron tokens. The token of unit $i$ in a fully connected
layer $\ell$ is its incoming weights and bias,
$[W^{(\ell)}_{i,:},\,b^{(\ell)}_i]$, and in a convolutional layer it is one
filter with its bias. Any vector $u$ over the weights is split in the same way,
with $[u]^{(\ell)}_i$ its token at unit $i$ of layer $\ell$. The forward drift is
predicted token by token,
\begin{equation}
    \big[v^{\mathrm{f}}_\phi(\Theta_s,s)\big]^{(\ell)}_i
    =f^{(\ell)}_\phi\!\left([\Theta_s]^{(\ell)}_i,\,s,\,[g_s]^{(\ell)}_i,\,
    [a_s]^{(\ell)}_i\right),
    \label{eq:tokendrift}
\end{equation}
where $g_s=\nabla_\theta\log p(\mathcal{D}_{t+1}\mid\Theta_s)$, $a_s$ is the
prior term of (\ref{eq:pathscore}) at $\Theta_s$, and $f^{(\ell)}_\phi$ is a
multilayer perceptron shared by the $n_\ell$ units of layer $\ell$, which also
receives $s$. Sharing turns each particle into $n_\ell$ training examples. The
score inputs tell the network how the new task and the carried posterior pull on
each unit, so the drift depends on this evidence rather than on a task
identifier and remains usable at the next boundary. A zero-initialized output
layer makes the untrained drift zero, so fitting starts from the reference
(\ref{eq:fisherreference}). The backward drift has the same form.

\subsection{Bridge-guided transitions}
\label{sec:transition}

At the first boundary, no drift has been learned and the transport is that of
Section~\ref{sec:langevin}. At every later boundary $t$, the particles are moved
by the Euler--Maruyama discretization of the forward drift fitted at boundary
$t-1$,
\begin{equation}
\begin{split}
    \Theta_{m}=\Theta_{m-1}
    &+v^{\mathrm{f}}_{\phi_{t-1}}(\Theta_{m-1},s_{m-1})\,\Delta s_m\\
    &+\sigma\sqrt{\Delta s_m}\,A_t^{-1/2}\xi,
\end{split}
    \label{eq:euler}
\end{equation}
with $s_m=\beta_m$ and $\Delta s_m=s_m-s_{m-1}$. Setting $s_m=\beta_m$ only
reuses the previous drift along the temperature schedule, and the weights
account for the difference between this learned transport and the tempered
posteriors,
\begin{equation}
    w^{(k)}\leftarrow w^{(k)}\,
    \frac{\pi_{t,\beta_m}(\Theta_m)\,L_m(\Theta_{m-1}\mid\Theta_m)}
         {\pi_{t,\beta_{m-1}}(\Theta_{m-1})\,K_m(\Theta_m\mid\Theta_{m-1})},
    \label{eq:bridgeweight}
\end{equation}
where $K_m$ and $L_m$ are the Gaussian transition densities of (\ref{eq:euler})
and of the corresponding step of the backward drift~\cite{delmoral2006smc}. With
these weights, the terminal population targets $\tilde q_{t+1}$ regardless of
the accuracy of the drift, while an accurate drift keeps the weights nearly
uniform, so that fewer Langevin steps and weight resets are needed. Weight
resets and Langevin steps (\ref{eq:mala}) follow each move as in
Section~\ref{sec:langevin}. After the transport, the drifts are refitted by IMF
on the new pair of populations, starting from their previous parameters. The
terminal particles with their normalized weights and the precision $A_{t+1}$
from (\ref{eq:precision}) form $\hat q_{t+1}$, after which $\mathcal{D}_{t+1}$ is
discarded. The final prediction is the weighted average of the predictive
distributions of the particles,
$p(y\mid x)\approx\sum_k w^{(k)}_{t+1}p(y\mid x,\theta^{(k)}_{t+1})$.
